\documentclass[preview]{standalone}

\usepackage{amsmath}
\usepackage{amssymb}
\usepackage{stellar}
\usepackage{definitions}
\usepackage{tikz}

\usetikzlibrary{cd}

\begin{document}

\id{tensor-products}
\genpage

\section{Tensor product}

\begin{snippetdefinition}{tensor-product-definition}{Tensor product}
    Let \(V_1, \cdots, V_n\) be \vectorspace[vector spaces]
    over a \field \(\mathbb F\). Then, their \emph{tensor product}
    is defined as a \vectorspace over \(\mathbb F\) 
    \[
        \mathbb F\text{-}\bigotimes_{i=1}^n V_i
    \]
    such that there exists an \(n\)-linear canonical map
    \[
        \varphi \colon \prod_{i=1}^n V_i \fromto \mathbb F\text{-}\bigotimes_{i=1}^n V_i
    \]
    and such that for every other \vectorspace \(W\) over \(\mathbb F\)
    and \(n\)-linear application \[
        h \colon \prod_{i=1}^n V_i \fromto W
    \]
    there exists a unique linear application \(\tilde h\) such that the following diagram commutes:
    \begin{center}
        % https://tikzcd.yichuanshen.de/#N4Igdg9gJgpgziAXAbVABwnAlgFyxMJZABgBpiBdUkANwEMAbAVxiRAB120AnaAfWBYAvAEYAvgD0wAAgBqfLCDGl0mXPkIoR5KrUYs2nAEZYA5hDwBbeAOHipchUpUgM2PASLaRu+s1aIIADqSrowUKbwRKAAZryWSGQgOBBI2nr+huz03GgAForKsfFp1ClIAEzUDFhgASBQdHB54SDUfgaBnHgMsNJ5zsUQCYhJ5YhVIAx0RjAMAApqHpog3GZ5OG0ZnSADYhRiQA
        \begin{tikzcd}
        \prod_{i=1}^n V_i \arrow[r, "\varphi"] \arrow[rd, "h"'] & \bigotimes_{i=1}^n V_i \arrow[d, "\tilde h", dashed] \\
                                                                & W                                                   
        \end{tikzcd}
    \end{center}
    The elements of \(\bigotimes_{i=1}^n V_i\) are called \emph{tensors}.
    Given \(v_1 \in V_1, \cdots, v_n \in V_n\), the \emph{simplest tensors}
    or \emph{pure tensors} are tensors of the form
    \[
        \bigotimes_{i=1}^n v_i
    \]
\end{snippetdefinition}

\begin{snippet}{tensor-product-intuition}
    The tensor product is conceptually the largest
    and most general vector space designed to linearize multilinearity.
    Its core purpose is to act as a bridge that transforms \(n\)-linear maps 
    into standard linear maps. By ``packaging'' the interactions between 
    vectors from different spaces into a single object, the tensor product 
    ensures that any multilinear relationship is captured within a purely 
    linear framework. It is considered the ``largest'' such space because it is 
    constructed without imposing any additional algebraic constraints or 
    relations beyond those strictly required by the laws of multilinearity 
    (such as distributivity and scalar displacement).
    In general, tensors have the form of a finite sum of pure tensors
\end{snippet}

\end{document}